\subsection{15.65: nonrational classical-group limits}
\label{cand:15.65}
For the full families \(U(n,q)\), \(\mathrm{Sp}(2m,q)\),
\(O(2m+1,q)\), \(O^+(2m,q)\), and \(O^-(2m,q)\), both limiting
proportions asked about are nonrational functions of \(q\).
For symplectic and orthogonal groups the assertion holds separately
on even and on odd prime powers.
The Notebook calls a matrix separable when its \emph{minimal}
polynomial is squarefree; this is semisimplicity.
The enumeration source uses separable for squarefree
\emph{characteristic} polynomial. We distinguish the two and prove
nonrationality for both, as well as for cyclic matrices.
This is a statement about functions of field size, not about
irrationality of an individual probability.

All enumeration and limit-existence formulas below are those of
Fulman--Neumann--Praeger (FNP) \cite{fnp}.
The proposed additional deduction is their analytic continuation
in inverse field size and the obstruction to rational interpolation.
Let \(t=1/q\). The count polynomials for unitary factors are
\[
\begin{aligned}
 N_d(q)&=\frac1d\sum_{a\mid d}\mu(a)(q^{d/a}+1)
                       &&(d\text{ odd}),\\
 M_1(q)&=(q^2-q-2)/2,\\
 M_d(q)&=\frac1{2d}\sum_{a\mid d}\mu(a)(q^{2d/a}-q^{d/a})
                       &&(d>1\text{ odd}),\\
 M_d(q)&=\frac1{2d}\sum_{a\mid d}\mu(a)q^{2d/a}
                       &&(d\text{ even}).
\end{aligned}
\]
For symplectic and orthogonal factors fix \(e=1\) for even \(q\)
or \(e=2\) for odd \(q\), and use
\[
\begin{aligned}
 J_d(q)&=\frac1{2d}\sum_{\substack{a\mid d\\a\text{ odd}}}
                  \mu(a)(q^{d/a}+1-e),\\
 K_1(q)&=(q-e-1)/2,\\
 K_d(q)&=\tfrac12 N(q;d)&& (d>1\text{ odd}),\\
 K_d(q)&=\tfrac12\bigl(N(q;d)-J_{d/2}(q)\bigr)&& (d\text{ even}),\\
 N(q;d)&=\frac1d\sum_{a\mid d}\mu(a)q^{d/a}&& (d>1).
\end{aligned}
\]
These are FNP Lemmas~1.3.12 and~1.3.16. In particular
\[
 N_1(1/t)=1+1/t,\qquad J_1(1/t)=(1/t+1-e)/2.
\]
For \(0<|t|<1\), their explicit formulas give bounds by a constant
times \(|t|^{-d}\) for \(N_d,J_d,K_d\), and \(|t|^{-2d}\) for \(M_d\).
For instance, bounding the number of divisors by \(d\) suffices.
These estimates will justify the infinite products as analytic functions.

\paragraph{The analytic obstruction.}
Suppose a holomorphic function \(a(t)\) has a zero of multiplicity
\(h>0\) at \(b\ne0\), \(P(t)\) is holomorphic and nonconstant near
\(b\), and \(K(t)\) is a nonzero single-valued meromorphic germ there.
Then a branch of \(K(t)\exp(P(t)\log a(t))\) cannot continue a rational
germ. One circuit around \(b\) multiplies it by
\(\exp(2\pi i hP(t))\). Rational single-valuedness would make this
identically one on an open set, forcing \(P'=0\).
An integer-order zero or pole of \(K\) cannot cancel this monodromy.

Each germ constructed below has a finite nonzero value at \(t=0\)
and agrees with the indicated probability at sufficiently small
positive \(t=1/q\). If a rational function of \(q\) agreed on all
powers of a fixed prime, its transform in \(t\) would have no pole
at zero and would agree with the germ by the identity theorem.
The analytic obstruction therefore excludes even that restricted
form of rational interpolation.

\paragraph{Cyclic limits.}
Put
\[
 A(u)=\frac{1+u-u^2}{1+u},\qquad
 B(u)=\frac{1-u+u^2}{1-u}=\frac{1+u^3}{1-u^2}.
\]
FNP Theorems~2.1.9 and~2.2.9 give
\[
\begin{aligned}
 C(U;q)&=(1+t)
   \prod_{\substack{d\ge1\\d\text{ odd}}}A(t^d)^{N_d(1/t)}
   \prod_{d\ge1}B(t^{2d})^{M_d(1/t)},\\
 C(\mathrm{Sp};q)&=
   \prod_{d\ge1}A(t^d)^{J_d(1/t)}B(t^d)^{K_d(1/t)}.
\end{aligned}
\]
Powers mean exponentials of logarithms initially vanishing at zero.
The first \(A\)-factor has a simple zero at
\(\alpha=(1-\sqrt5)/2\).
Remove it. On \(|t|<3/4\), every remaining \(A\)-factor has
degree \(d\ge2\), and
\[
 \left|\frac{t^{2d}}{1+t^d}\right|
 \le\frac{(3/4)^4}{1-(3/4)^2}=\frac{81}{112}<1.
\]
The logarithms of all \(B\)-factors are single-valued on \(|t|<1\)
by the second expression for \(B\).
Their logarithms vanish to order at least \(2d\), or \(4d\) after
substitution \(t^{2d}\). Multiplication by the count polynomials
therefore leaves summable bounds by constant multiples of
\(|t|^d\) or \(|t|^{2d}\), uniformly on compact subdiscs.
All apparent singularities at zero are removable.
The remaining product is consequently holomorphic and nonvanishing
on \(|t|<3/4\). The first exponents \(N_1(1/t)\) and \(J_1(1/t)\)
are nonconstant, so the obstruction at \(\alpha\) proves
nonrationality. In the unitary case the real local exponent is
\(1+1/\alpha=\alpha\), also giving a direct noninteger-order proof.
FNP Theorem~2.3.11 gives
\[
 C(O_{\mathrm{odd}};q)=(1-t)C(\mathrm{Sp};q),\qquad
 C(O^\pm;q)=\frac{(1-t)^e+1}{2}C(\mathrm{Sp};q).
\]
These nonzero rational multipliers transfer the conclusion to all
full orthogonal families.

\paragraph{A zero needed for the literal separability question.}
Define holomorphic functions on \(|u|<1\) by
\[
 \mathcal U(u)=\sum_{m\ge0}
       \frac{u^{m^2}}{\prod_{i=1}^m(1-(-u)^i)},\qquad
 \mathcal L(u)=\sum_{m\ge0}
       \frac{u^{m^2}}{\prod_{i=1}^m(1-u^i)}.
\]
On \(|u|\le r<1\), denominator products are bounded below by
\(\prod_{i\ge1}(1-r^i)>0\), proving normal convergence.
Set
\(a_m(r)=r^{m^2}/\prod_{i=1}^m(1-r^i)\).
For \(m\ge1\),
\[
 \frac{a_{m+1}(r)}{a_m(r)}
 \le\frac{r^3}{1-r^2},\qquad
 \sum_{m\ge1}a_m(3/7)\le\frac{210}{253}<1.
\]
Hence both functions are nonzero on \(|u|\le3/7\).
On the negative real axis,
\(\mathcal U(-r)=1-a_1(r)+a_2(r)-\cdots\).
For \(r\le3/5\) these positive terms decrease. Alternating bounds
give positivity for \(r\le1/2\), a lower bound \(9/56\) at \(1/2\),
and the following negative upper bound at \(3/5\):
\[
 1-a_1+a_2-a_3+a_4
 =-\frac{1128797279}{26656000000}<0.
\]
Thus there is a first negative zero \(b\in(-3/5,-1/2)\), of finite
positive multiplicity, with \(\mathcal U(t)>0\) for \(b<t\le0\).
Also \(\mathcal L(-r)>1-r/(1+r)>0\) for \(0<r\le13/20\):
its alternating term ratios are bounded by \(r^3/(1-r^2)<1\).
No simplicity or exact location of \(b\) is needed.

Put
\[
 E(u)=(1-u)\mathcal U(u),\qquad
 D(u)=(1-u)\mathcal L(u).
\]
They are nonzero on \(|u|\le3/7\), and both are \(1+O(u^2)\).
Let \(SS\) denote semisimplicity, namely squarefree minimal polynomial.
FNP Theorems~3.1.13 and~3.1.15 give
\[
\begin{aligned}
 SS(U;q)&=(1+t)
  \prod_{\substack{d\ge1\\d\text{ odd}}}E(t^d)^{N_d(1/t)}
  \prod_{d\ge1}D(t^{2d})^{M_d(1/t)},\\
 SS(\mathrm{Sp};q)&=(1-t)^e\mathcal P(t)^e T_e(t),\\
 T_e(t)&=\prod_{d\ge1}E(t^d)^{J_d(1/t)}D(t^d)^{K_d(1/t)},\\
 \mathcal P(t)&=1+\sum_{m\ge1}
       \frac{t^{m(2m+1)}}{\prod_{i=1}^m(1-t^{2i})}.
\end{aligned}
\]
The function \(\mathcal P\) is single-valued holomorphic on the unit
disc, with value one at zero.
Take \(R=13/20\). All arguments of factors of degree at least two
lie in \(|u|\le R^2=169/400<3/7\).
Their logarithms and the count-polynomial bounds thus give normally
convergent, nonvanishing remainders on \(|t|<R\).
For \(T_e\), separate also the degree-one \(D(t)\)-factor.
It has a holomorphic logarithm near \([b,0]\), by the positivity
just proved; multiplication by \(K_1(1/t)\) is removable at zero
since \(\log D(t)=O(t^2)\).
The only required branching at \(b\) comes from \(E(t)\), with
nonconstant exponent \(N_1(1/t)\) or \(J_1(1/t)\).
The prefactors can have only integer-order zeros there.
The obstruction proves both nonrationality claims.

For the orthogonal case put
\[
 \mathcal Q(t)=1+\sum_{m\ge1}
       \frac{t^{m(2m-1)}}{\prod_{i=1}^m(1-t^{2i})}.
\]
It is likewise holomorphic on the unit disc, with value one at zero.
The reciprocal orders of the two even-dimensional full orthogonal
groups sum to the terms of \(\mathcal Q(1/q)-1\).
FNP Theorem~3.1.18 gives
\[
\begin{aligned}
 SS(O_{\mathrm{odd}};q)&=
   (1-t)^e\mathcal Q(t)^{e-1}\mathcal P(t)T_e(t),\\
 SS(O^\pm;q)&=
   \frac{(1-t)^e}{2}
   \bigl(\mathcal Q(t)^e+(e-1)\mathcal P(t)^2\bigr)T_e(t).
\end{aligned}
\]
The extra factors are single-valued and not identically zero:
at zero they equal \(1\) and \(e/2\).
Their possible finite-order zeros cannot cancel the monodromy of
\(E(t)^{J_1(1/t)}\). This proves the literal Notebook assertion for
all the stated families.

\paragraph{The other meaning of separability.}
For completeness let \(RS\) mean squarefree characteristic polynomial,
and set \(V(u)=(1-u)(1+2u)/(1+u)\).
FNP Theorems~2.1.3, 2.2.3 and~2.3.4 give
\[
\begin{aligned}
 RS(U;q)&=(1+t)
       \prod_{\substack{d\ge1\\d\text{ odd}}}V(t^d)^{N_d(1/t)},\\
 RS(\mathrm{Sp};q)&=(1-t)^e
       \prod_{d\ge1}V(t^d)^{J_d(1/t)},\\
 RS(O_{\mathrm{odd}};q)&=RS(\mathrm{Sp};q),\\
 RS(O^\pm;q)&=2^{e-2}RS(\mathrm{Sp};q).
\end{aligned}
\]
The first \(V\)-factor vanishes at \(-1/2\).
All higher factors have zero-free logarithms on \(|t|<13/20\),
and the same estimates prove normal convergence.
The first exponents are nonconstant, so the branching obstruction
applies even when their values at \(-1/2\) happen to be integers.
This proves nonrationality for the source's convention too.

The symplectic semisimple formula uses \(M^*(q;d)\), as in FNP
Theorem~3.1.15 and its underlying product; an intermediate display
in Lemma~3.1.9(e) prints \(M^*(q;2d)\).
The archive records that index discrepancy and a separate coefficient
discrepancy. No numerical coefficient table is a premise here.
The distinction between the dimension variable of FNP's generating
functions and our field-size variable \(t\) is essential.
See \artifact{research/15.65-proof.md},
\artifact{research/15.65-classical-proof.md}, and their source audits.
